% TOP_PROBLEM: {"id": 4800035, "title": "Multiple ergodic averages — Problem 35", "category_id": 11, "category": {"id": 11, "name": "dynamical_systems", "display_name": "Dynamical Systems", "description": "Problems about long-term behavior of deterministic systems, Hamiltonian dynamics, and periodic orbits.", "slug": "dynamical-systems", "order_index": 11, "created_at": "2026-07-31T15:26:25.671Z"}, "status": "open", "problem_number": "AMR-047-0035", "set_id": 13, "statusQualification": "Restores the omitted set A after the middle transformation in the source display; positive indices suffice, so the unused T0 notation is omitted."}
% FIRST_POSTED: 2026-10-01
% ENTRY_KIND: statement_only
% UnsolvedMath ID 4800035; source catalogue code AMR-047-0035
% !TeX program = lualatex
% Definition and sources only; this file is not an LLM solution attempt.
\documentclass[11pt,a4paper]{article}
\usepackage[margin=25mm]{geometry}
\usepackage{fontspec}
\setmainfont{lmroman10-regular.otf}[BoldFont=lmroman10-bold.otf,
 ItalicFont=lmroman10-italic.otf,BoldItalicFont=lmroman10-bolditalic.otf]
\usepackage{amsmath,amssymb,mathtools}
\usepackage{unicode-math}
\setmathfont{latinmodern-math.otf}
\AtBeginDocument{\let\setminus\smallsetminus}
\usepackage{xurl,hyperref}
\hypersetup{colorlinks=true,urlcolor=blue}
\setcounter{secnumdepth}{0}
\setlength{\parindent}{0pt}
\setlength{\parskip}{5pt}
\setlength{\emergencystretch}{3em}
\begin{document}
\section*{Multiple ergodic averages — Problem 35}\label{4800035}
{\small UnsolvedMath ID 4800035 · AMR-047-0035 · Dynamical Systems · AMR Open Problem Lists}
\subsection{Definitions and mathematical statement}
A probability-preserving system with multiplicative
structure consists of a probability space
$(X,\mathcal B,\mu)$ and invertible
measure-preserving transformations $T_r$, one
for every positive integer $r$, such that
\[
                  T_1=\mathrm{id},\qquad
                  T_rT_s=T_{rs}\quad(r,s\ge1).
\]
Thus all the transformations commute. Equivalently,
they form an action of the multiplicative positive
integers by invertible maps; this extends to
positive rational indices using inverses.

For every such system and every measurable
$A$ with $\mu(A)>0$, must there exist
positive integers $m,n$ such that
\[
 \mu\left(
 T_{m(m+n)}A\ \cap\ T_{(m+2n)(m+3n)}A\ \cap\
 T_{(m+4n)(m+5n)}A
 \right)>0?
\]
All three expressions denote images of the
same set under the corresponding indexed
transformation. The missing $A$ after the
middle transformation in the catalogue/source
display is supplied explicitly.

The indices are products of consecutive terms
of a single six-term arithmetic progression;
they are not additive iterates $T^r$ of
one fixed map. Both $m$ and $n$ must
be positive, so setting $n=0$ cannot give
a vacuous affirmative answer. No ergodicity,
mixing, or freeness of the multiplicative
action is required. The question asks only
for positive measure of one three-way
intersection, without prescribing a lower
bound, density of good pairs, or a limit
for their averages.
\subsection{Short English statement}
Must multiplicative actions give a positive three-way intersection at the three paired-product indices of a six-term arithmetic progression?
\subsection{Sources}
\begin{itemize}
\item[S1] ULAM.AI and the UnsolvedMath Contributors, supplied problems.json, record 4800035 (AMR-047-0035), AMR Open Problem Lists. Catalogue identity and source transcription; CC BY 4.0.
\url{https://huggingface.co/datasets/ulamai/UnsolvedMath}.
\item[S2] Frantzikinakis - Some open problems on multiple ergodic averages (2016), Problem 35. Original source indexed by AMR-047-0035.
\url{https://arxiv.org/abs/1103.3808}.
\end{itemize}
\end{document}
